\documentclass[../../main]{subfiles}

\begin{document}

\section{Famílias de Funções}
\label{sec:matematica:funcoes:familias-de-funcoes}

\subsection{Famílias Indexadas}

\begin{defn}[Família de Funções]
	\label{def:matematica:funcoes:familias:familia}

	Sejam \(I\), \(A\) e \(B\) conjuntos.

	Uma família de funções de \(A\) em \(B\)
	indexada por \(I\) é uma função

	\[
		F:I\to B^A.
	\]

	Para cada \(i\in I\), escreve-se

	\[
		f_i = F(i).
	\]

	A família é então denotada por

	\[
		(f_i)_{i\in I}.
	\]

\end{defn}

\begin{obs}

	Cada elemento do conjunto índice determina uma função da família.

\end{obs}

\begin{ex}

	Seja

	\[
		f_n:\mathbb R\to\mathbb R,
		\qquad
		f_n(x)=x^n,
	\]

	para cada

	\[
		n\in\mathbb N.
	\]

	Então

	\[
		(f_n)_{n\in\mathbb N}
	\]

	é uma família de funções indexada pelos números naturais.

\end{ex}

\subsection{Famílias Finitas}

\begin{defn}[Família Finita]
	\label{def:matematica:funcoes:familias:finita}

	Uma família

	\[
		(f_i)_{i\in I}
	\]

	é dita finita quando o conjunto índice \(I\) é finito.

\end{defn}

\begin{ex}

	As funções

	\[
		f_1,f_2,f_3
	\]

	formam uma família finita indexada por

	\[
		I=\{1,2,3\}.
	\]

\end{ex}

\subsection{Sequências de Funções}

\begin{defn}[Sequência de Funções]
	\label{def:matematica:funcoes:familias:sequencia}

	Uma sequência de funções é uma família indexada por

	\[
		\mathbb N.
	\]

	Isto é,

	\[
		(f_n)_{n\in\mathbb N}.
	\]

\end{defn}

\begin{obs}

	Sequências de funções desempenham papel fundamental em análise
	matemática, especialmente nos estudos de convergência pontual e
	convergência uniforme.

\end{obs}

\subsection{Famílias Constantes}

\begin{defn}[Família Constante]
	\label{def:matematica:funcoes:familias:constante}

	Seja

	\[
		f:A\to B.
	\]

	A família

	\[
		(f_i)_{i\in I}
	\]

	é dita constante quando

	\[
		f_i=f
	\]

	para todo

	\[
		i\in I.
	\]

\end{defn}

\begin{ex}

	Seja

	\[
		f(x)=0.
	\]

	A família

	\[
		(f_n)_{n\in\mathbb N}
	\]

	definida por

	\[
		f_n=f
	\]

	para todo \(n\) é constante.

\end{ex}

\subsection{Igualdade de Famílias}

\begin{defn}[Igualdade de Famílias]
	\label{def:matematica:funcoes:familias:igualdade}

	Sejam

	\[
		(f_i)_{i\in I}
	\]

	e

	\[
		(g_i)_{i\in I}
	\]

	famílias de funções.

	Diz-se que essas famílias são iguais quando

	\[
		f_i=g_i
	\]

	para todo

	\[
		i\in I.
	\]

\end{defn}

\begin{prop}
	\label{prop:matematica:funcoes:familias:extensionalidade}

	Se

	\[
		f_i(x)
		=
		g_i(x)
	\]

	para todo

	\[
		i\in I
	\]

	e para todo

	\[
		x\in A,
	\]

	então

	\[
		(f_i)_{i\in I}
		=
		(g_i)_{i\in I}.
	\]

\end{prop}

\begin{proof}

	Para cada \(i\in I\),

	\[
		f_i(x)
		=
		g_i(x)
	\]

	para todo \(x\in A\).

	Pelo princípio da extensionalidade das funções,

	\[
		f_i=g_i.
	\]

	Como isso vale para todo índice,

	\[
		(f_i)_{i\in I}
		=
		(g_i)_{i\in I}.
	\]

\end{proof}

\subsection{Famílias de Subconjuntos}

\begin{obs}

	Uma família de subconjuntos de um conjunto \(A\) é um caso
	particular de família indexada:

	\[
		F:I\to\mathcal P(A).
	\]

	Essa observação será útil posteriormente em Topologia,
	Álgebra e Teoria da Medida.

\end{obs}

\end{document}
